\section*{Chapitre \ref{ch:s2:harvesting-the-variance-risk-premium} --- Récolter la prime de risque de variance}

\begin{solution}{exo:s2:harvesting-the-variance-risk-premium:1}
$0.8 \times 0.20 \times \sqrt{21/252} = 0.8 \times 0.20 \times 0.289 = 4.6~\%$ de l'indice.
\end{solution}

\begin{solution}{exo:s2:harvesting-the-variance-risk-premium:2}
$0.16^2 \times 1.3 = 0.0333$, une volatilité implicite de $\sqrt{0.0333} = 18.2~\%$ contre une valeur attendue de 16~\%.
\end{solution}

\begin{solution}{exo:s2:harvesting-the-variance-risk-premium:3}
L'option de vente est hors de la monnaie~: elle expire sans valeur dans la plupart des mois (92~\%) et encaisse la prime entière, mais, lors d'un krach, sa perte est plusieurs fois supérieure à la prime, une
asymétrie de $-7.2$~; le P\&L du swap de variance est plus régulier d'un mois à l'autre, si bien que sa moyenne représente une plus grande part de sa volatilité.
\end{solution}

\begin{solution}{exo:s2:harvesting-the-variance-risk-premium:4}
Un écart-type mensuel vaut $10~\%/\sqrt{12} = 2.9~\%$~; 7.7 d'entre eux représentent 22.2~\% du capital.
\end{solution}

\begin{solution}{exo:s2:harvesting-the-variance-risk-premium:5}
Les krachs simulés partent de marchés calmes, où la variance implicite et donc le strike sont faibles, si bien que le portefeuille est le plus important précisément au moment où le krach
débute. Un dimensionnement inversement proportionnel au strike aiderait si de fortes pertes suivaient une forte volatilité, comme lorsqu'une période agitée se prolonge plutôt que
qu'un calme prenne fin.
\end{solution}

\begin{solution}{exo:s2:harvesting-the-variance-risk-premium:6}
Les investisseurs paient pour se protéger contre la variance, qui monte lorsqu'ils sont les plus pauvres. La moyenne du VIX moins la volatilité réalisée inclut les mois où
la volatilité réalisée a largement dépassé le VIX~; la prime se mesure sur la variance, où ces mois pèsent davantage, et sur de longues périodes qui incluent des krachs.
\end{solution}

\begin{solution}{exo:s2:harvesting-the-variance-risk-premium:7}
À 30~\% le pire mois coûte 66.5~\% et le capital survit~; à 50~\% il coûte 110.8~\% et le capital est perdu. Le pire mois vaut 7.68 écarts-types
mensuels, de sorte que la totalité du capital disparaît pour une cible de $1/(7.68/\sqrt{12}) = 45~\%$.
\end{solution}

\begin{solution}{exo:s2:harvesting-the-variance-risk-premium:8}
Cinq années calmes ne disent rien du krach que la stratégie est rémunérée pour supporter~; l'historique synthétique du portefeuille, avant son pire mois, présentait aussi un
ratio de Sharpe supérieur à 1.5. Il faut dimensionner le fonds d'après un stress (un mois comme celui de février 2020 ou d'octobre 2008 appliqué aux positions actuelles), et non d'après son historique.
\end{solution}

\begin{solution}{pb:s2:harvesting-the-variance-risk-premium:1}
\begin{enumerate}
\item Vendre des options ou de la variance pour gagner la prime~; vendre des options contre une position détenue, le plus souvent des calls.
\item La prime de risque de variance, définie comme la variance réalisée moins un taux de swap de variance synthétisé à partir d'options, sur cinq indices et 35 actions.
\item VIX à 19.5 contre une volatilité réalisée de 15.4 en moyenne, supérieur dans 84~\% des mois, de 4.1 points~; variance supérieure de 31~\%.
\item février 2020, septembre et août 2008, mars 2025, juillet 2011, juin 2002.
\item Variance stochastique avec effet de levier, sauts et trois krachs~; variance implicite égale à 1.3 fois la variance attendue, un smile de volatilité décroissant.
\item 0.69, 1.24, 0.40, 0.08 et 0.09~; asymétrie allant de $-2.3$ à $-12.5$.
\item La prime rémunère la prise en charge des krachs~: de petits gains la plupart des mois, de fortes pertes, rarement.
\item PutWrite à 7.2~\% par an (ratio de Sharpe 0.57), BuyWrite à 6.0~\% (0.48)~; $-27.6~\%$ et $-28.5~\%$ à l'automne 2008, $-19.8~\%$ et $-21.3~\%$ au début de 2020.
\item Le rendement d'une option couverte en delta, rapporté à sa prime, qui mesure la prime de volatilité.
\item Les options longues sur indice couvertes en delta sous-performent zéro, moins lorsqu'elles sont hors de la monnaie, davantage lorsque la volatilité est élevée.
\item 0.53~\% de l'indice par mois, soit environ un huitième de la prime~; la couverture supprime la direction de l'indice.
\item Moins de transactions de couverture et davantage de dépendance au chemin~: coûts plus faibles, P\&L plus bruité.
\item \textbf{Résultat nommé.} Le ratio de Sharpe du portefeuille vendeur de variance, avant son pire mois, vaut 1.57 quelle que soit sa taille~; le pire mois coûte 11.1~\%, 22.2~\%,
44.3~\% et 88.6~\% du capital pour des volatilités de 5~\%, 10~\%, 20~\% et 40~\%, et plus que le capital à 60~\%.
\item Le krach est un événement unique et rare~; cinq années calmes ne contiennent aucune information sur son ampleur.
\item Par un stress appliqué aux positions actuelles, de sorte que le pire mois plausible soit supportable.
\item Il a augmenté sa taille sur les marchés calmes, là où les krachs ont commencé.
\item Les années de krach de l'échantillon, des prix d'options exécutables, des appels de marge et une comparaison directe avec le sous-jacent.
\item La vente de puts l'emporte le plus souvent (92~\% des mois)~; le straddle couvert présente le meilleur ratio de Sharpe.
\item Le ciblage de volatilité réduit le risque après une hausse de la volatilité~; ici les pertes surviennent pendant que la volatilité monte, plus vite qu'une cible ne peut réagir.
\item Une assurance contre la variance.
\end{enumerate}
\end{solution}

\begin{solution}{iq:s2:harvesting-the-variance-risk-premium:1}
La différence entre le prix de la variance (implicite) et la variance réalisée ensuite~; elle est positive en moyenne parce que la variance est élevée dans les mauvais moments
et les investisseurs paient pour s'en couvrir.
\end{solution}

\begin{solution}{iq:s2:harvesting-the-variance-risk-premium:2}
La perte de gamma du straddle vaut environ $\tfrac12\Gamma(\Delta S)^2$, importante pour un écart de 5~\%, et la volatilité implicite monte aussi (une perte de vega)~; il faut recouvrir le nouveau
delta à l'ouverture, vérifier la limite de stress et réduire le portefeuille s'il la dépasse, plutôt que d'attendre un rebond.
\end{solution}

\begin{solution}{iq:s2:harvesting-the-variance-risk-premium:3}
Avec des limites de stress sur la perte lors de krachs historiques et hypothétiques appliquées aux positions actuelles, des limites de vega et de gamma par maturité, un plafond sur la
prime encaissée par rapport au capital, et des limites de liquidité sur la vitesse à laquelle le portefeuille peut être réduit.
\end{solution}

\begin{solution}{iq:s2:harvesting-the-variance-risk-premium:4}
Les swaps de variance captent la prime sur l'ensemble des strikes et ont un profil de gain net et convexe, avec une exposition au krach dans la queue de distribution~; les straddles la captent à la monnaie
avec couverture en delta et dépendance au chemin~; les options de vente hors de la monnaie captent la partie la plus riche du skew, avec le risque de krach le plus concentré.
\end{solution}

\begin{solution}{iq:s2:harvesting-the-variance-risk-premium:5}
Positions et prix, une surface par sous-jacent, grecques par option, agrégation par sous-jacent, une règle de couverture (temporelle ou par bandes), génération d'ordres dans
l'instrument de couverture, et rapprochement des exécutions de couverture avec le modèle.
\end{solution}

\begin{solution}{iq:s2:harvesting-the-variance-risk-premium:6}
La variation de la position couverte vaut $\tfrac12\Gamma(dS)^2 + \Theta\,dt$~; avec $\Theta = -\tfrac12\Gamma S^2\sigma_i^2$ (Black--Scholes à volatilité implicite) et
$(dS)^2 = S^2\sigma_r^2\,dt$, le P\&L vaut $\tfrac12\Gamma S^2(\sigma_r^2 - \sigma_i^2)\,dt$~: une option longue gagne lorsque la volatilité réalisée dépasse la volatilité implicite, une option courte
gagne lorsque la volatilité implicite dépasse la volatilité réalisée, pondéré par le gamma.
\end{solution}
